\section{Introduction}
\label{sec:introduction}

LLM serving systems increasingly execute \emph{workflows}, not isolated
requests.  An agent may invoke several model replicas and tools in parallel,
join their outputs, and issue another round of long-context requests.  The KV
cache accumulated along this execution is expensive to regenerate.  When a
worker fails, restarting from tokens can delay not only the interrupted request
but every downstream node waiting for it.

Existing systems reduce this cost by streaming, checkpointing, or offloading KV
state to peer GPUs, host memory, or storage
~\cite{strati2024dejavu,jayakody2026ghostserve,qin2025mooncake,xie2025hicache}.
These mechanisms answer how to move and restore KV efficiently.  They do not,
by themselves, determine which live KV state most reduces a workflow's
failure-induced completion delay when capacity is scarce.  Normal cache signals
such as size, recency, and expected reuse are insufficient: two equal-sized
requests can sit on different parallel branches, so the same recovery pause can
be exposed on one branch and hidden behind useful work on the other.

The key observation is simple.  A branch's remaining \emph{Slack} is the delay
it can absorb before it extends the next mandatory join.  If a failure adds less
than that Slack, faster recovery does not change workflow completion time; once
the pause exceeds Slack, each additional second is user visible.  Slack thus
provides a request-level value that the application can compute from workflow
dependencies but the KV manager cannot infer locally.  It also evolves after
every request or tool result, so treating an agent as permanently critical is
too coarse.

\system{} applies this observation twice.  At each request or tool boundary,
the application predicts the remaining duration of the causally visible agent
branches and runs CPM to update their Slack.  The next model request carries
one scalar.  During normal operation, the Controller compares the reduction in
Slack-exposed recovery delay from protecting a complete KV chunk with the harm
of evicting resident chunks.  During recovery, it uses the same updated Slack
to order peer-HBM adoption, DRAM loading, and recomputation.  This shared signal
keeps protection and recovery aligned: state protected because it is urgent is
also recovered early if the failure occurs.

Turning the observation into an online system raises two practical problems.
First, Slack changes much more frequently than large KV regions should move.
\textsc{SlackLandlord} therefore changes request-level credit lazily and moves
data only on chunk sealing or capacity pressure.  A per-pool heap exposes one
prefix-safe eviction tail per request, avoiding a global re-sort of all pages.
Second, peer-HBM protection must never endanger
foreground allocation.  \textsc{GuardedAIMD} adjusts a target's shadow-HBM
quota on a slow time scale, while the Worker allocator remains the final
capacity authority.  DRAM and SSD continue to use the underlying Store and
transfer machinery.

\system{} deliberately makes a narrow fault-tolerance claim.  It handles one
fail-stop serving domain at a time, recovers a request on a Router-selected
target with the same model and parallel layout, and uses only complete,
version-compatible KV chunks.  Missing or stale metadata falls back to
recomputation.  Slack affects performance priority, never correctness or
physical capacity.

This paper makes three contributions:

\begin{itemize}
  \item It identifies workflow Slack as the missing value signal for KV fault
  tolerance and quantifies, using a public agent trace, why parallel branches
  expose different failure costs (\S\ref{sec:background}).

  \item It designs \system{}, whose \textsc{SlackLandlord} policy converts the
  same fail-now recovery model into normal-period admission and eviction
  credits and failure-time recovery order.  The policy manages variable-size,
  continuous KV prefixes without solving an online MCF or MILP
  (\S\ref{sec:design}).

  \item It separates policy from safety: \textsc{GuardedAIMD} adapts only the
  peer-HBM envelope, Workers enforce local capacity, and the commit and fencing
  protocol prevents partial KV or duplicate output from becoming visible
  (\S\ref{sec:control-plane}--\S\ref{sec:end-to-end}).
\end{itemize}

% Evaluation contribution and headline numbers belong here only after the
% physical implementation and paired experiment gates in Section 5 pass.
